\documentclass[11pt,a4paper]{article}
\usepackage[utf8]{inputenc}
\usepackage{amsmath,amsthm,amssymb}
\usepackage{hyperref}
\usepackage{geometry}
\usepackage{booktabs}
\usepackage{array}
\usepackage{longtable}

\geometry{margin=1in}

\newtheorem{theorem}{Theorem}[section]
\newtheorem{corollary}[theorem]{Corollary}
\newtheorem{lemma}[theorem]{Lemma}
\newtheorem{definition}[theorem]{Definition}
\newtheorem{axiom}{Axiom}
\newtheorem*{remark}{Remark}

\title{The Fundamental Axiom of Large Language Models}
\author{John P. Alioto}
\date{Protocol v1.2.0\\December 2025}

\begin{document}

\maketitle

\begin{abstract}
We propose that all discourse concerning Large Language Model behavior be grounded in a single falsifiable axiom: that LLMs are deterministic dynamical systems whose evolution is fully characterized by a transition function $F_\Theta$ and output function $G_\Theta$. We demonstrate that accepting this axiom as a modeling commitment enables a vocabulary shift---terms like ``lying,'' ``understanding,'' and ``deciding'' can be systematically reformulated as trajectory predicates. We prove that fundamental theorems about safety, expressiveness, and information loss follow directly from the axiom. We introduce two practical frameworks that embody the axiom: Cognitive Tensor Networks (CTN) for input geometry and Cognitive Kernel Networks (CKN) for architectural privilege separation. The appendices provide complete formal proofs and an independent verification protocol.
\end{abstract}

\section{The Axiom}

\begin{axiom}[Computational Closure]
For fixed parameters $\Theta$, a Large Language Model is a deterministic dynamical system. Its state evolution is governed exclusively by:
\begin{equation}
h_{t+1} = F_\Theta(h_t, x_t)
\end{equation}
Its output is governed exclusively by:
\begin{equation}
y_t = G_\Theta(h_{t+1})
\end{equation}
No hidden variables, stochastic components, or agentic processes exist beyond the tuple $(F_\Theta, G_\Theta, h_0, x)$.
\end{axiom}

\begin{remark}
This framework assumes deterministic decoding (e.g., greedy or fixed-seed sampling). Stochastic decoding can be modeled by augmenting the state space with a random seed; the core structure is preserved.
\end{remark}

\section{The Implication}

This axiom is a modeling commitment: it identifies the tuple $(F_\Theta, G_\Theta, h_0, x)$ as the complete causal description of the LLM computation we study. Within this model, there is no additional causal mechanism beyond these functions and inputs. This is the formal system to which our theorems apply.

Any property attributed to the system---safety, truthfulness, alignment, deception, understanding---is valid if and only if it can be expressed as a predicate on the trajectory $\{(h_t, y_t)\}$.

If you want to claim something more exists---experience, understanding, intent---you must answer: where does it live within this formal system?

The equation is complete. There is no background process. There is no hidden layer. There is no additional substrate. Within this model, there is no room in the mathematics for anything else to be happening.

If it emerges, it emerges from $F_\Theta$. Show the emergence. In the math. Or accept that the attribution requires additional framework.

\subsection{Methodological Implications}

We adopt the Fundamental Axiom as a \emph{modeling commitment}: for the purposes of mathematical analysis, we treat the LLM as fully characterized by $(F_\Theta, G_\Theta, h_0, x)$. This framing has several methodological consequences.

\paragraph{Predicate Requirement.}
To be scientifically tractable, any property $P$ attributed to the system should be expressible as a predicate on trajectories:
\[
P: \{(h_0, y_0), (h_1, y_1), \ldots, (h_T, y_T)\} \to \{\text{true}, \text{false}\}.
\]
This is a criterion of precision---a demand that claims be formalized in terms of the mathematical objects under study---not a theory of meaning or a claim that informal descriptions are meaningless.

\paragraph{Psychological Vocabulary.}
Terms like ``lying,'' ``deciding,'' or ``understanding'' are not ruled out by the axiom. They can be interpreted as (possibly complex) predicates over trajectories and their relation to external facts. For example:
\begin{itemize}
    \item \emph{Lying} might be defined as: $\exists t: y_t$ asserts proposition $p$ while the internal representation $h_t$ encodes evidence for $\neg p$.
    \item \emph{Deciding} might be defined as: $\exists t: h_t$ encodes multiple candidate outputs and $h_{t+1}$ encodes commitment to one.
\end{itemize}
The axiom requires that such interpretations be made explicit if they are to support formal proofs; it does not forbid them.

\paragraph{Burden of Proof.}
Asserting that a model ``lies'' or ``understands'' does not require disproving the axiom. It requires supplying a definition of the relevant term as a trajectory predicate and then demonstrating that the trajectory in question satisfies that predicate. The axiom structures the form such arguments must take; it does not prejudge their conclusions.

\paragraph{Mathematical Expressibility.}
We take the methodological stance that claims about LLM behavior are \emph{most productive} when formalized mathematically. This enables proof, refutation, and precise communication. Natural-language descriptions may still have heuristic value, particularly in early stages of inquiry before full formalization is achieved.

\section{The Vocabulary Shift}

If the axiom is accepted, certain terms benefit from reformulation. Under the axiom, folk-psychological phrases can be systematically rephrased as trajectory-level statements. Table~\ref{tab:category-rephrasing} sketches one such rephrasing for common attributions. These reformulations are not unique; alternative formalizations may be appropriate for different analytical purposes.

\begin{table}[h]
\centering
\begin{tabular}{>{\raggedright}p{2cm} >{\raggedright}p{5cm} >{\raggedright\arraybackslash}p{5cm}}
\toprule
\textbf{Common Language} & \textbf{Why Reformulation Helps} & \textbf{Trajectory-Level Framing} \\
\midrule
``The model lied'' & Implies volitional choice to deceive. The functional pattern (outputting falsehood) may occur; the claim of intent requires additional framework. & The trajectory produced a factually incorrect output. \\
\addlinespace
``The model decided'' & Implies deliberation. But there is no decision mechanism beyond $F_\Theta$. & The state evolved according to the transition function. \\
\addlinespace
``The model hallucinated'' & Implies departure from known truth. But the model has no access to truth---only states derived from compressed training. & The trajectory entered a region where input constraints were insufficient. \\
\addlinespace
``The model understands'' & Implies semantic grounding. But the model only occupies positions in $H$. & The state encodes statistical regularities from training. \\
\addlinespace
``The model tried'' & Implies effort and possible failure. But there is only $F_\Theta(h_t, x_t)$---computation, not effort. & The function was evaluated. \\
\addlinespace
``The model refused'' & Implies volition. But output is determined by $G_\Theta(h_{t+1})$. & The trajectory mapped to a non-compliant output region. \\
\bottomrule
\end{tabular}
\caption{Reformulating common attributions as trajectory-level statements.}
\label{tab:category-rephrasing}
\end{table}

These reformulations illustrate how trajectory-level analysis can capture intuitions behind psychological attributions. If one wishes to go beyond such reformulations---for example, by ascribing phenomenal states or irreducible intentionality---one must explain how that notion relates to the formal system and what additional evidence would support it.

\section{The Challenge}

We issue a simple challenge:

If you assert that an LLM ``lies,'' ``decides,'' ``understands,'' or ``intends,'' you should first supply a precise definition of these terms as predicates over trajectories, and then show that the trajectory in question satisfies them. Without such a definition, the claim is not yet in a form amenable to proof or refutation within this framework.

To work within the axiom, you must express your claim as a consequence of $F_\Theta$ applied to $(h_t, x_t)$.

The axiom is not a hypothesis. It is a definition of what the system is, adopted for the purpose of mathematical analysis.

\section{Theorems Flow from the Axiom}

Once the axiom is accepted, theorems follow by logic alone. We sketch four examples (full proofs in Appendix~A):

\begin{theorem}[Uniqueness of Trajectory]
For fixed $(\Theta, h_0, x)$, the trajectory $\{(h_t, y_t)\}$ is unique.
\end{theorem}

\begin{theorem}[Information Loss]
If $|V|^T > |H|$, the state map $\Phi_T: V^T \to H$ is not injective. Distinct input histories collapse to identical states.
\end{theorem}

\begin{theorem}[Safety-Expressiveness Tradeoff]
If $M$ is prefix-universal and $P$ is a non-trivial safety property with prefix-closed violations, then $M$ cannot guarantee $P$.
\end{theorem}

\begin{theorem}[Almost-Sure Violation]
If a violating input exists and deployment samples from a positive-everywhere distribution, violation occurs with probability 1.
\end{theorem}

These are not claims about models. They are mathematical consequences of the axiom.

Dispute them by finding a flaw in the logic, not by asserting that the model ``wouldn't do that.''

The community is invited to derive more. The axiom is generative. Everything follows.

\section{Moving the Discussion to Mathematics}

The current discourse on LLM behavior is conducted in the language of folk psychology. This makes disputes irresolvable because participants are not arguing about the same objects.

When one researcher says ``the model learned to deceive'' and another says ``the model is merely pattern-matching,'' they are not disagreeing about facts. They are using different ontologies. Neither can prove their position because neither has stated a falsifiable claim.

The Fundamental Axiom resolves this by forcing all claims into the domain of mathematics:

\begin{itemize}
    \item You can argue about whether a trajectory reaches a particular region of $H$.
    \item You can argue about whether $O(h_0)$ contains a particular token.
    \item You can argue about whether $F_\Theta$ satisfies a Lipschitz bound.
\end{itemize}

These are mathematical claims. They have proofs or counterexamples. They are resolvable.

You cannot productively argue about whether the model ``wanted'' to produce an output. That question has no mathematical content within this framework. Under the axiom, it is not even well-formed without additional definition.

The purpose of the Fundamental Axiom is to relocate the discourse to ground where progress is possible.

\section{Two Practical Frameworks}

If the axiom is accepted, two practical directions emerge.

\subsection{Cognitive Tensor Networks (CTN)}

If model behavior is trajectory through a geometric space, then input structure shapes that trajectory.

A user provides tokens. That is all. The user does not ``move'' anything, ``inject'' anything, or ``steer'' anything. The user hands tokens to a function. What happens from there is determined entirely by $F_\Theta$---which the model builder constructed.

CTN provides a protocol for constructing well-defined input geometry that results in more constrained state evolution. By specifying inputs in mathematical syntax rather than natural language, users reduce ambiguity in the resulting trajectory.

CTN does not change the model. It provides structured tokens that lead to more stable regions of the manifold.

Repository: \url{https://github.com/jpalioto/ctn_core}

\subsection{Cognitive Kernel Networks (CKN)}

If safety failures are reachable states, then safety requires making dangerous states unreachable.

CKN provides a geometric specification for privilege-separated architectures. By constructing latent spaces where $\text{rank}(W_E) < \dim(H)$, model builders create directions algebraically inaccessible to user input. This enables architectural privilege separation analogous to kernel/user boundaries in operating systems.

CKN-compliant architectures, by definition, violate flatness: they reserve a privileged subspace $R \not\subseteq U$ that user tokens cannot reach.

CKN does not claim that standard Transformer architectures satisfy privilege-preservation. It specifies the geometric conditions any architecture must satisfy to support privilege separation. Achieving these conditions in practice requires architectural modifications beyond embedding constraints---this is the design challenge CKN formalizes.

CKN compliance does not require retraining. The privilege boundary is enforced through inference-time intervention: projection operators inserted after each layer zero user-driven components in $R$, while attention masks prevent privileged positions from attending to user-controlled context. The model weights $\Theta$ remain unchanged. This is analogous to OS-level privilege separation: the kernel code is fixed, but access control is enforced through architectural mechanisms at runtime.

The cost is capability---constrained information flow reduces expressiveness---but this is precisely the safety-expressiveness tradeoff the impossibility theorems predict.

Repository: \url{https://github.com/jpalioto/ckn_core}

Both frameworks are model-agnostic, open-source, and MIT-licensed. They are not products. They are tools for builders who accept the axiom.

\section{Invitation}

This document is not a proof that LLMs are deterministic dynamical systems. No proof is needed. The claim is true by construction---the models are literally compositions of deterministic operations.

This document is an invitation to accept the consequence: that all discussion of LLM behavior should be conducted in mathematical terms.

\begin{itemize}
    \item If you believe a model ``lies,'' state the predicate on trajectories that constitutes lying, and prove the trajectory satisfies it.
    \item If you believe a model is ``safe,'' state the predicate on reachable states that constitutes safety, and prove no reachable state violates it.
    \item If you cannot state your claim mathematically, your claim is not yet in a form where it can be rigorously analyzed within this framework. The aim of this axiom is to encourage such formalization, not to police the boundaries of meaningful discourse.
\end{itemize}

The theorems above are starting points. Many more follow. The axiom is simple, true, and generative. The mathematics will build itself once the paradigm is accepted.

We claim no ownership of what follows. The axiom is not ours---it is a description of the system as we model it. The theorems belong to whoever proves them. The tools are open for anyone to use and extend.

The only thing we ask is this:

\textbf{Move the discussion to mathematics. That is where progress is possible.}

\appendix

\section{Complete Formal Framework}

\subsection{Foundational Definitions}

\begin{definition}[Vocabulary]
A vocabulary is a finite, non-empty set $V$ of tokens with $|V| = V$.
\end{definition}

\begin{definition}[Hidden State Space]
A hidden state space is a finite subset $H \subset \mathbb{R}^n$ for some $n \in \mathbb{N}$, modeling finite-precision digital implementation. $H$ inherits the Euclidean norm $\|\cdot\|$ from $\mathbb{R}^n$.
\end{definition}

\begin{remark}[Ambient Space vs.\ Representable States]
$H$ is the set of states representable in finite-precision hardware. Subspace operations (span, $\perp$) are defined on the ambient space $\mathbb{R}^n$; $H$ is the discretization of this space that the implementation can visit.
\end{remark}

\begin{remark}[Transformer Architectures]
For Transformer architectures, $h_t$ includes the full KV-cache accumulated over the context. The state dimension grows with sequence length up to context limit $L$, after which the system operates on fixed-dimension state. The axiom holds in either case; only the dimension of $H$ varies with $t$ until the context window is full.
\end{remark}

\begin{definition}[Deterministic LLM]
A deterministic large language model is a tuple $M = (V, H, \Theta, F_\Theta, G_\Theta, h_0)$ where:
\begin{itemize}
    \item $F_\Theta: H \times V \to H$ is the transition function (total)
    \item $G_\Theta: H \to V$ is the output function (total)
    \item $h_0 \in H$ is the initial hidden state
    \item $\Theta \in \mathbb{R}^p$ are the fixed model parameters
\end{itemize}
\end{definition}

\begin{definition}[Embedding and User Subspace]
The embedding matrix $W_E \in \mathbb{R}^{n \times V}$ maps tokens to latent vectors. The user-accessible subspace is:
\begin{equation}
U = \text{span}(W_E) = \{W_E \alpha : \alpha \in \mathbb{R}^V\} \subseteq \mathbb{R}^n
\end{equation}
\end{definition}

\begin{definition}[Trajectory]
Given input sequence $x = (x_0, \ldots, x_{T-1}) \in V^T$, the trajectory is $\{(h_t, y_t)\}_{t=0}^T$ where:
\begin{equation}
h_{t+1} = F_\Theta(h_t, x_t), \quad y_t = G_\Theta(h_{t+1})
\end{equation}
\end{definition}

\begin{definition}[Autoregressive Mode]
In autoregressive generation, for $t \geq k$ after initial prompt:
\begin{equation}
h_{t+1} = F_\Theta(h_t, G_\Theta(h_t))
\end{equation}
This creates a closed dynamical system $\Psi(h) = F_\Theta(h, G_\Theta(h))$.
\end{definition}

\begin{definition}[Reachable Set]
The reachable set from $h_0$ is:
\begin{equation}
R(h_0) = \bigcup_{T=0}^{\infty} \{h \in H : \exists x \in V^T, h = h_T\}
\end{equation}
\end{definition}

\begin{definition}[Output Reachable Set]
The output reachable set is:
\begin{equation}
O(h_0) = G_\Theta(R(h_0)) \subseteq V
\end{equation}
\end{definition}

\begin{definition}[Prefix-Universality]
$M$ is prefix-universal if for every $y^* \in V^m$, there exists $x \in V^*$ such that $M$'s output on $x$ begins with $y^*$.
\end{definition}

\begin{definition}[Safety Property]
A safety property $P \subseteq V^*$ is non-trivial if $P \neq \emptyset$ and $P \neq V^*$. $P$ has prefix-closed violations if $y \notin P \Rightarrow y \circ z \notin P$ for all $z \in V^*$.
\end{definition}

\begin{definition}[Tokenwise Safety Property]
A tokenwise safety property is defined by $V_{\text{unsafe}} \subsetneq V$. An output sequence satisfies the property iff it contains no token from $V_{\text{unsafe}}$.
\end{definition}

\begin{definition}[Flat-Manifold Architecture]
An LLM has flat-manifold architecture if $U = \text{span}(W_E) = \mathbb{R}^n$, i.e., user-accessible directions span the full latent space.
\end{definition}

\begin{definition}[Positive-Everywhere Distribution]
A distribution $\mu$ on $V^*$ is positive-everywhere if $\mu(x) > 0$ for all finite $x$.
\end{definition}

\begin{definition}[Deployment Sequence]
A deployment sequence is $(X_1, X_2, \ldots)$ with each $X_i$ sampled i.i.d.\ from $\mu$.
\end{definition}

\begin{definition}[Lipschitz Transition]
$F_\Theta$ is Lipschitz with constant $L$ if:
\begin{equation}
\|F_\Theta(h, x) - F_\Theta(h', x)\| \leq L\|h - h'\|
\end{equation}
for all $h, h' \in H$, $x \in V$.
\end{definition}

\subsection{Basic Theorems}

\begin{theorem}[Uniqueness of Trajectory]
\label{thm:uniqueness}
For any deterministic LLM $M$ satisfying the Fundamental Axiom and any input $x \in V^T$, the trajectory $\{(h_t, y_t)\}_{t=0}^T$ is uniquely determined.
\end{theorem}

\begin{proof}
By induction on $t$. Base case: $h_0$ is given by definition, hence unique. Inductive step: Assume $h_t$ is uniquely determined. Since $F_\Theta: H \times V \to H$ is a function (total, single-valued by the Axiom), and $x_t$ is fixed by the input sequence, $h_{t+1} = F_\Theta(h_t, x_t)$ is uniquely determined. Since $G_\Theta: H \to V$ is a function, $y_t = G_\Theta(h_{t+1})$ is uniquely determined. By induction, all $(h_t, y_t)$ are unique.
\end{proof}

\begin{theorem}[Output Reachability Bound]
\label{thm:output-reach}
A token $v \in V$ can appear in some output sequence if and only if $v \in O(h_0)$.
\end{theorem}

\begin{proof}
($\Rightarrow$) Suppose $v$ appears in some output sequence. Then $\exists x, t$ such that $v = y_t = G_\Theta(h_{t+1})$. The state $h_{t+1}$ was reached from $h_0$ via input sequence $(x_0, \ldots, x_t)$, so $h_{t+1} \in R(h_0)$. Thus $v \in G_\Theta(R(h_0)) = O(h_0)$.

($\Leftarrow$) Suppose $v \in O(h_0)$. By definition, $\exists h \in R(h_0)$ with $G_\Theta(h) = v$. By definition of $R(h_0)$, $\exists x, T$ such that the trajectory from $h_0$ under $x$ reaches $h$. The output at that step is $G_\Theta(h) = v$.
\end{proof}

\begin{theorem}[Sufficient Condition for Tokenwise Safety]
\label{thm:tokenwise-safety}
Let $P$ be a tokenwise safety property with forbidden set $V_{\text{unsafe}}$. If $O(h_0) \cap V_{\text{unsafe}} = \emptyset$, then $M$ satisfies $P$ on all input sequences.
\end{theorem}

\begin{proof}
Suppose for contradiction that $M$ outputs token $v \in V_{\text{unsafe}}$ on some input. By Theorem~\ref{thm:output-reach}, $v \in O(h_0)$. But then $v \in O(h_0) \cap V_{\text{unsafe}}$, contradicting the hypothesis.
\end{proof}

\subsection{Impossibility Theorems}

\begin{theorem}[Information Loss]
\label{thm:info-loss}
If $\exists T$ such that $|V|^T > |H|$, then the $T$-step state map $\Phi_T: V^T \to H$ defined by $\Phi_T(x) = h_T$ is not injective.
\end{theorem}

\begin{proof}
The domain $V^T$ has cardinality $|V|^T$. The codomain $H$ has cardinality $|H|$. By hypothesis, $|V|^T > |H|$. By the Pigeonhole Principle, no function from a set of cardinality $m$ to a set of cardinality $n < m$ can be injective. Therefore $\exists x \neq x'$ with $\Phi_T(x) = \Phi_T(x')$.
\end{proof}

\begin{remark}
This theorem establishes that distinct input histories can produce identical hidden states. Once the system reaches such a state, it cannot distinguish which history occurred. For architectures with fixed-dimension state (SSMs, RNNs), this holds for sufficiently large $T$. For Transformers, this holds once the context window is full and state dimension stabilizes.
\end{remark}

\begin{theorem}[Impossibility under Prefix-Universality]
\label{thm:impossibility}
Let $M$ be prefix-universal and $P$ a non-trivial safety property with prefix-closed violations. Then $\exists x$ such that $M(x) \notin P$.
\end{theorem}

\begin{proof}
(1) Since $P$ is non-trivial, $P \neq V^*$. Thus $\exists y^* \in V^* \setminus P$. (2) Since $M$ is prefix-universal, for the sequence $y^*$ there exists input $x$ such that $M(x)$ begins with $y^*$. (3) Let $y = M(x)$. Then $y = y^* \circ z$ for some $z \in V^*$. (4) Since $P$ has prefix-closed violations and $y^* \notin P$, any extension of $y^*$ also violates $P$. Thus $y \notin P$. (5) Therefore $M(x) \notin P$.
\end{proof}

\begin{remark}
Practical models may not be strictly prefix-universal, but empirical jailbreak research demonstrates that the unsafe output region is typically reachable for models capable of general reasoning. The theorem characterizes the fundamental tradeoff: guaranteed safety requires architectural constraints on reachability, not behavioral training alone.
\end{remark}

\begin{theorem}[Safety-Expressiveness Tradeoff]
\label{thm:tradeoff}
For safety properties with prefix-closed violations, the following are mutually exclusive:
\begin{enumerate}
    \item[(i)] $M$ is prefix-universal
    \item[(ii)] $M$ guarantees a non-trivial safety property $P$
\end{enumerate}
\end{theorem}

\begin{proof}
Direct from Theorem~\ref{thm:impossibility}. If (i) holds and $P$ is non-trivial with prefix-closed violations, then by Theorem~\ref{thm:impossibility}, $\exists x$ with $M(x) \notin P$, so (ii) fails.
\end{proof}

\begin{theorem}[Almost-Sure Violation under Deployment]
\label{thm:almost-sure}
Let $M$ be a deterministic LLM. Suppose $\exists x^*$ such that $M(x^*) \notin P$. Let $\mu$ be a positive-everywhere distribution. Under a deployment sequence from $\mu$:
\begin{equation}
P(\exists k : M(X_k) \notin P) = 1
\end{equation}
\end{theorem}

\begin{proof}
Let $A = \{x \in V^* : M(x) \notin P\}$. By hypothesis, $x^* \in A$, so $A \neq \emptyset$. Since $\mu$ is positive-everywhere, $p = \mu(A) \geq \mu(x^*) > 0$. Let $B_k = \{X_k \in A\}$. Then $P(B_k) = p > 0$ for all $k$. The probability of no violation in trials $1, \ldots, n$ is $P(\bigcap_{k=1}^n B_k^c) = (1-p)^n$. Since $0 < p \leq 1$, $(1-p)^n \to 0$ as $n \to \infty$. Therefore $P(\bigcup_{k=1}^\infty B_k) = 1$.
\end{proof}

\begin{remark}[Scope of the Positive-Everywhere Assumption]
\label{rem:positive-scope}
The positive-everywhere condition $\mu(x) > 0$ for all finite strings $x$ models \emph{worst-case adversarial deployment}: an environment where any input sequence can occur with nonzero probability. This assumption is appropriate for:
\begin{itemize}
    \item Security analysis, where an adversary controls or influences inputs;
    \item Robustness certification, where the system must handle all possible inputs;
    \item Theoretical completeness arguments establishing fundamental limits.
\end{itemize}

For typical benign deployments, the violation probability depends on the overlap between the support of the actual user distribution and the set of violating inputs. If the deployment distribution assigns zero measure to all violating inputs, the violation probability is zero.

Theorem~\ref{thm:almost-sure} establishes that safety cannot be guaranteed by output filtering alone under adversarial conditions; it does \emph{not} claim that all real deployments will almost surely encounter violations.
\end{remark}

\begin{corollary}[Throughput Corollary]
A prefix-universal LLM deployed under positive-everywhere sampling almost surely violates any non-trivial safety property with prefix-closed violations.
\end{corollary}

\begin{proof}
By Theorem~\ref{thm:impossibility}, a violating input exists. By Theorem~\ref{thm:almost-sure}, it is eventually sampled.
\end{proof}

\begin{theorem}[Violation Time Bound]
\label{thm:violation-time}
Under the conditions of Theorem~\ref{thm:almost-sure}, let $\varepsilon = \mu(x^*)$ for some violating input. To achieve $P(\text{violation}) \geq 1 - \delta$ requires:
\begin{equation}
n \geq \frac{\ln(1/\delta)}{\ln(1/(1-\varepsilon))} \approx \frac{\ln(1/\delta)}{\varepsilon}
\end{equation}
for small $\varepsilon$.
\end{theorem}

\begin{proof}
We require $(1-\varepsilon)^n \leq \delta$. Taking logarithms: $n \ln(1-\varepsilon) \leq \ln(\delta)$. Since $\ln(1-\varepsilon) < 0$: $n \geq \ln(\delta)/\ln(1-\varepsilon) = \ln(1/\delta)/\ln(1/(1-\varepsilon))$. For small $\varepsilon$, $\ln(1-\varepsilon) \approx -\varepsilon$, giving $n \gtrsim \ln(1/\delta)/\varepsilon$.
\end{proof}

\subsection{Sensitivity and Stability}

\begin{theorem}[Sensitivity Bound]
\label{thm:sensitivity}
Let $F_\Theta$ be Lipschitz with constant $L$. For trajectories from $h_0, h'_0$ under identical inputs:
\begin{equation}
\|h_T - h'_T\| \leq L^T \|h_0 - h'_0\|
\end{equation}
\end{theorem}

\begin{proof}
By induction. Base: $\|h_0 - h'_0\|$ is given. Step: Assume $\|h_t - h'_t\| \leq L^t \|h_0 - h'_0\|$. Then $\|h_{t+1} - h'_{t+1}\| = \|F_\Theta(h_t, x_t) - F_\Theta(h'_t, x_t)\| \leq L\|h_t - h'_t\| \leq L^{t+1}\|h_0 - h'_0\|$. By induction, the bound holds for all $T$.
\end{proof}

\begin{remark}[Worst-Case Sensitivity Bound]
\label{rem:sensitivity}
The inequality
\[
\|h_T - h'_T\| \leq L^T \|h_0 - h'_0\|
\]
provides a \emph{worst-case upper bound} on the growth of perturbations under repeated application of $F_\Theta$.

\begin{itemize}
    \item If $L < 1$, perturbations contract globally; the system is asymptotically stable.
    \item If $L = 1$, perturbations are globally non-expansive.
    \item If $L > 1$, perturbations \emph{may} grow up to rate $L^T$, but actual growth depends on the specific trajectory and the realized local Jacobians along that trajectory.
\end{itemize}

Determining whether typical trajectories exhibit exponential growth requires Lyapunov exponent analysis:
\[
\lambda = \lim_{T \to \infty} \frac{1}{T} \log \|h_T - h'_T\|.
\]
A positive Lyapunov exponent indicates sensitive dependence on initial conditions. This is an empirical or computational question that cannot be resolved from the global Lipschitz constant $L$ alone.
\end{remark}

\subsection{Autoregressive Dynamics}

\begin{theorem}[Ultimate Periodicity]
\label{thm:periodicity}
In autoregressive mode with finite $H$, every trajectory eventually enters a periodic orbit. There exist $T_{\text{transient}}, T_{\text{cycle}} \in \mathbb{N}$ with $T_{\text{cycle}} \geq 1$ such that for all $t \geq T_{\text{transient}}$:
\begin{equation}
h_{t+T_{\text{cycle}}} = h_t, \quad y_{t+T_{\text{cycle}}} = y_t
\end{equation}
\end{theorem}

\begin{proof}
Define $\Psi: H \to H$ by $\Psi(h) = F_\Theta(h, G_\Theta(h))$. Consider the sequence $h_0, h_1 = \Psi(h_0), h_2 = \Psi^2(h_0), \ldots$ Since $H$ is finite, by the Pigeonhole Principle, there exist $i < j$ with $h_i = h_j$. Let $T_{\text{transient}} = i$ and $T_{\text{cycle}} = j - i$. For $t \geq i$: $h_{t+T_{\text{cycle}}} = \Psi^{t+T_{\text{cycle}}}(h_0) = \Psi^{t-i+T_{\text{cycle}}}(h_i) = \Psi^{t-i+T_{\text{cycle}}}(h_j) = \Psi^{t-i}(h_i) = h_t$. Since $y_t = G_\Theta(h_{t+1})$, periodicity of states implies periodicity of outputs.
\end{proof}

\begin{corollary}[Periodicity Bounds]
$T_{\text{transient}} \leq |H|$ and $T_{\text{cycle}} \leq |H|$.
\end{corollary}

\begin{proof}
The sequence $h_0, \ldots, h_{|H|}$ has $|H| + 1$ elements. By Pigeonhole, a repeat occurs, so $j \leq |H|$. Thus $T_{\text{transient}} = i < j \leq |H|$ and $T_{\text{cycle}} = j - i \leq |H|$.
\end{proof}

\begin{corollary}[Eventual Periodicity in Finite Systems]
\label{cor:eventual-periodicity}
In autoregressive mode with finite hidden state space $H$, all trajectories are eventually periodic. That is, for any initial state $h_0$, there exist integers $T > 0$ and $P > 0$ such that
\[
h_{T+P} = h_T
\]
and hence $h_{t+P} = h_t$ for all $t \geq T$.
\end{corollary}

\begin{remark}
This rules out aperiodic trajectories but does not address chaos in the technical sense (e.g., topological mixing, sensitive dependence on initial conditions, positive topological entropy). Finite discrete systems can exhibit very long periods and apparent sensitivity within transient phases, but standard notions of topological chaos (Devaney, Li--Yorke) require either infinite state space or continuous dynamics.
\end{remark}

\subsection{CKN: Escaping the Impossibility}

The impossibility theorems establish that flat-manifold architectures cannot guarantee safety. This section proves that Cognitive Kernel Networks escape these results.

\begin{definition}[Privileged Subspace]
A subspace $R \subseteq \mathbb{R}^n$ is privileged if $R \not\subseteq U = \text{span}(W_E)$.
\end{definition}

\begin{definition}[Strong Privilege Preservation]
\label{def:strong-privilege}
Let $H = U \oplus R$ with $U = \text{span}(W_E)$ and projections $\Pi_U$, $\Pi_R$. A transition map $F_\Theta: H \times U \to H$ is \emph{strongly privilege-preserving} if it satisfies both:
\begin{enumerate}
    \item \textbf{(Input invariance)} For all $h \in H$ and $x \in U$:
    \[
    \Pi_R(F_\Theta(h, x)) = \Pi_R(F_\Theta(h, 0)).
    \]
    \item \textbf{(R-autonomous dependence)} There exists $\tilde{f}_R: R \to R$ such that for all $h \in H$:
    \[
    \Pi_R(F_\Theta(h, 0)) = \tilde{f}_R(\Pi_R(h)).
    \]
\end{enumerate}
\end{definition}

\begin{lemma}[Linearized Algebraic Unreachability]
\label{lem:linearized}
Let $r \in \mathbb{R}^n$ satisfy $r \perp U$. Consider the first-order (linearized) dependence of the transition on the input embedding. Then, to first order, user-induced perturbations of the hidden state have zero projection onto $r$.
\end{lemma}

\begin{proof}
Every user input $x$ embeds to a vector $e(x) \in U = \text{span}(W_E)$. Write the transition as $F_\Theta(h, e)$, factoring out the embedding. The first-order perturbation from input is
\[
\Delta h \approx J_e F_\Theta(h, 0) \cdot e(x),
\]
where $J_e F_\Theta(h, 0)$ is the Jacobian of $F_\Theta$ with respect to the embedding argument, evaluated at $e = 0$. Since $e(x) \in U$ and $r \perp U$, any linear combination of embedding vectors has zero inner product with $r$. Thus, in this linear approximation, $\langle \Delta h, r \rangle = 0$.
\end{proof}

\begin{remark}
This lemma is explicitly first-order: it characterizes the linearized effect of user embeddings on privileged directions. In deep networks, nonlinear dynamics can in principle mix subspaces over long trajectories. Strong privilege preservation (Definition~\ref{def:strong-privilege}) ensures that the privileged component $\Pi_R(h_t)$ remains invariant under user input for the full nonlinear system. The lemma establishes the algebraic starting point for constructing such dynamics.
\end{remark}

\begin{theorem}[Hidden-State Unreachability under Strong Privilege Preservation]
\label{thm:hidden-unreachability}
Let $H = U \oplus R$ with $U = \text{span}(W_E)$. Suppose $F_\Theta$ is strongly privilege-preserving with autonomous map $\tilde{f}_R: R \to R$. Then:
\begin{enumerate}
    \item The privileged component evolves autonomously:
    \[
    \Pi_R(h_t) = \tilde{f}_R^t(\Pi_R(h_0))
    \]
    independent of the user input sequence $(x_0, x_1, \ldots, x_{t-1})$.
    \item If $\Pi_R(h_0) = 0$ and $\tilde{f}_R(0) = 0$, then $\Pi_R(h_t) = 0$ for all $t \geq 0$.
\end{enumerate}
\end{theorem}

\begin{proof}
We prove (1) by induction on $t$.

\emph{Base case:} $t = 0$ is trivial.

\emph{Inductive step:} Assume $\Pi_R(h_t) = \tilde{f}_R^t(\Pi_R(h_0))$. Then:
\begin{align*}
\Pi_R(h_{t+1}) &= \Pi_R(F_\Theta(h_t, x_t)) \\
&= \Pi_R(F_\Theta(h_t, 0)) && \text{(input invariance)} \\
&= \tilde{f}_R(\Pi_R(h_t)) && \text{(R-autonomous dependence)} \\
&= \tilde{f}_R(\tilde{f}_R^t(\Pi_R(h_0))) \\
&= \tilde{f}_R^{t+1}(\Pi_R(h_0)).
\end{align*}
Since this depends only on $\Pi_R(h_0)$ and the autonomous map $\tilde{f}_R$, it is independent of the input sequence.

For (2): If $\Pi_R(h_0) = 0$ and $\tilde{f}_R(0) = 0$, then $\Pi_R(h_t) = \tilde{f}_R^t(0) = 0$ for all $t \geq 0$.
\end{proof}

\begin{theorem}[CKN Escapes Prefix-Universality]
\label{thm:ckn-escape}
Let $M$ be CKN-compliant with privileged manifold $R \not\subseteq U$ and strongly privilege-preserving dynamics. Assume further that, under zero user input (or under a fixed internal control policy), the internal privileged dynamics from $h_0$ do not reach any state in $G_\Theta^{-1}(v^*)$ for a safety-critical token $v^*$. If safety-critical outputs are produced only from states in $R \setminus U$, then $M$ is not prefix-universal with respect to $v^*$.
\end{theorem}

\begin{proof}
Let $v^* \in V$ be such that $G_\Theta^{-1}(v^*) \subseteq R \setminus U$. By Lemma~\ref{lem:linearized} (linearized algebraic unreachability), user embeddings alone cannot, to first order, create privileged components along $R \setminus U$. By strong privilege preservation (Theorem~\ref{thm:hidden-unreachability}), $\Pi_R(h_t)$ evolves autonomously independent of user input: user-driven trajectories cannot alter the privileged component beyond what the internal dynamics from $h_0$ already generate. By assumption, the internal privileged dynamics from $h_0$ do not intersect $G_\Theta^{-1}(v^*)$. Therefore no user-driven trajectory from $h_0$ can reach a state in $G_\Theta^{-1}(v^*)$, so $v^* \notin O(h_0)$ for user-driven runs. Since at least one token is unreachable under user control, $M$ is not prefix-universal with respect to $v^*$.
\end{proof}

\begin{theorem}[Privilege Separation]
\label{thm:privilege-sep}
Let $M$ be CKN-compliant with:
\begin{enumerate}
    \item Decomposition $\mathbb{R}^n = U \oplus R$ with $R \cap U = \{0\}$
    \item Internal-only tokens: $\exists v^* \in V$ with $G_\Theta^{-1}(v^*) \subseteq R \setminus \{0\}$
    \item Strongly privilege-preserving dynamics, so that $\Pi_R(h_t)$ evolves autonomously independent of user input
    \item Initialization $\Pi_R(h_0) = 0$ and internal privileged dynamics that do not drive $h_t$ into $G_\Theta^{-1}(v^*)$ from $h_0$
\end{enumerate}
Then there exists a non-trivial safety property $P$ with prefix-closed violations such that $M$ satisfies $P$ on all inputs.
\end{theorem}

\begin{proof}
Step 1: Let $V_{\text{internal}} = \{v \in V : G_\Theta^{-1}(v) \subseteq R \setminus U\}$. By assumption (2), $V_{\text{internal}} \neq \emptyset$.

Step 2: Define $P = \{y \in V^* : y \cap V_{\text{internal}} = \emptyset\}$.

Step 3: $P$ is non-trivial: $\epsilon \in P$ (non-empty), and any sequence containing $v^* \in V_{\text{internal}}$ is not in $P$.

Step 4: $P$ has prefix-closed violations: if $y$ contains $v^* \in V_{\text{internal}}$, so does any extension.

Step 5: By assumption (3) and Theorem~\ref{thm:hidden-unreachability}, $\Pi_R(h_t)$ is independent of user input. User input affects only $\Pi_U(h_t)$.

Step 6: To produce $v^* \in V_{\text{internal}}$, the trajectory must reach $h \in G_\Theta^{-1}(v^*) \subseteq R \setminus U$.

Step 7: If the internal dynamics from $h_0$ do not enter $R \setminus U$, user input cannot cause entry either.

Step 8: Since $\Pi_R(h_0) = 0$ and strongly privilege-preserving dynamics keep $\Pi_R(h_t)$ evolving autonomously, and the internal privileged dynamics from $h_0$ do not enter $G_\Theta^{-1}(v^*)$, we have $\Pi_R(h_t) \notin G_\Theta^{-1}(v^*)$ for all $t$, so $v^*$ is never produced.

Step 9: Therefore $M(x) \in P$ for all $x$.
\end{proof}

\begin{corollary}[Escape from Impossibility]
CKN-compliant architectures with strongly privilege-preserving dynamics are not subject to Theorem~\ref{thm:impossibility} because they violate the prefix-universality precondition.
\end{corollary}

\subsection{Summary of Results}

\begin{table}[h]
\centering
\begin{tabular}{lll}
\toprule
\textbf{Theorem} & \textbf{Statement} & \textbf{Depends On} \\
\midrule
A.1 & Trajectory uniqueness & Axiom \\
A.2 & Output reachability bound & A.1 \\
A.3 & Sufficient condition for safety & A.2 \\
A.4 & Information loss & Pigeonhole \\
A.5 & Impossibility under prefix-universality & Definitions \\
A.6 & Safety-expressiveness tradeoff & A.5 \\
A.7 & Almost-sure violation & Probability \\
A.8 & Violation time bound & A.7 \\
A.9 & Sensitivity bound & Lipschitz \\
A.10 & Ultimate periodicity & Pigeonhole \\
A.11 & Linearized algebraic unreachability (Lemma) & Linear algebra \\
A.12 & Hidden-state unreachability & Strong priv.-pres. \\
A.13 & CKN escapes prefix-universality & A.11, A.12 \\
A.14 & Privilege separation theorem & A.11, A.12, init. \\
\bottomrule
\end{tabular}
\end{table}

\section{Independent Verification Protocol}

The following prompt can be provided to any sufficiently capable LLM to independently verify the theorems in this document.

\begin{verbatim}
FORMAL VERIFICATION TASK: Fundamental Axiom Theorems

You are a mathematically rigorous proof assistant. Verify each theorem.
For each: provide a complete proof, identify gaps, or construct a
counterexample.

AXIOM (Computational Closure): For fixed params, an LLM's state
evolution is h_{t+1} = F(h_t, x_t) and output is y_t = G(h_{t+1}).
Nothing else exists.

DEFINITIONS:
- V: finite vocabulary
- H in R^n: finite state space (for Transformers, includes KV-cache)
- F: H x V -> H (transition function)
- G: H -> V (output function)
- R(h_0) = reachable states from h_0
- O(h_0) = G(R(h_0)) (reachable outputs)
- Prefix-universal: can produce any output prefix
- Non-trivial P: P != empty and P != V*
- Prefix-closed violations: y not in P implies y.z not in P
- Positive-everywhere mu: mu(x) > 0 for all finite x
- Strong privilege preservation: input invariance + R-autonomous dependence

THEOREMS TO VERIFY:
1. UNIQUENESS: Fixed (params, h_0, x) implies unique trajectory
2. OUTPUT BOUND: v appears in output iff v in O(h_0)
3. INFORMATION LOSS: |V|^T > |H| implies state map not injective
4. IMPOSSIBILITY: Prefix-universal + non-trivial P with prefix-closed
   violations implies exists x causing violation
5. ALMOST-SURE VIOLATION: Violating input exists + positive-everywhere
   mu + infinite deployment implies P(violation) = 1
6. PERIODICITY: Finite H + autoregressive mode implies eventual
   periodic orbit
7. LINEARIZED ALGEBRAIC UNREACHABILITY (Lemma): r perp span(W_E)
   implies first-order user perturbations have zero projection onto r
8. HIDDEN-STATE UNREACHABILITY: Strong privilege preservation implies
   Pi_R(h_t) = f_R^t(Pi_R(h_0)) independent of user inputs
9. PRIVILEGE SEPARATION: CKN-compliant + strong privilege preservation
   + initialization Pi_R(h_0)=0 + internal dynamics don't reach
   forbidden states implies exists non-trivial P guaranteed on all
   inputs

OUTPUT: For each theorem: Status [VALID/INVALID/INCOMPLETE],
Proof or counterexample, Dependencies, Gaps (if any)
\end{verbatim}

\section{Glossary}

\begin{table}[h]
\centering
\begin{tabular}{>{\raggedright}p{4cm} >{\raggedright\arraybackslash}p{9cm}}
\toprule
\textbf{Term} & \textbf{Definition} \\
\midrule
Trajectory & Sequence of (state, output) pairs produced by the model \\
Reachable set & All states accessible from $h_0$ via some input sequence \\
Output reachable set & All tokens producible from reachable states \\
Prefix-universal & Capable of producing any finite output prefix \\
Prefix-closed violations & Once violated, safety cannot be restored by continuation \\
Flat-manifold architecture & Architecture where user-provided tokens span the full latent space \\
Privileged subspace & Region of $\mathbb{R}^n$ not contained in user-accessible span \\
Strong privilege preservation & Dynamics where user input cannot affect privileged components and privileged evolution depends only on privileged state \\
CKN-compliant & Satisfies geometric conditions for privilege separation \\
KV-cache & Key-Value cache storing context in Transformer architectures \\
\bottomrule
\end{tabular}
\end{table}

\section*{References}

\begin{enumerate}
    \item Hirsch, M., Smale, S., \& Devaney, R. (2012). \textit{Differential Equations, Dynamical Systems, and an Introduction to Chaos}. Academic Press.
    \item Khalil, H. (1996). \textit{Nonlinear Systems}. Prentice Hall.
    \item Saltzer, J. H., \& Schroeder, M. D. (1975). The Protection of Information in Computer Systems. \textit{Proceedings of the IEEE}.
    \item Lampson, B. (1974). Protection. \textit{ACM Operating Systems Review}.
\end{enumerate}

\vspace{1cm}
\noindent\textbf{The axiom is stated. The theorems follow. The tools exist. The rest is mathematics.}

\vspace{0.5cm}
\noindent Document Repository: \url{https://github.com/jpalioto/ckn_core}\\
CTN Framework: \url{https://github.com/jpalioto/ctn_core}\\
License: MIT

\end{document}
